\documentclass[10pt]{beamer}

\usetheme{Madrid}
\usecolortheme{default}
\usefonttheme{professionalfonts}

\usepackage{amsmath,amssymb,amsfonts}
\usepackage{graphicx}
\usepackage{booktabs}
\usepackage{hyperref}

\title[Operator Learning]{Coding tutorial: Neural network optimization for PDEs}

\author[X. Xu]{Xiaofeng Xu\\
	~\\
%	{\scriptsize Postdoctoral Scholar, Department of Mathematics, Penn State University
%	}
%	\\
%	~\\
}

\institute[PSU]{Postdoctoral Scholar,\\
Department of Mathematics, Penn State University\\
%\textcolor{blue}{\tt xiaofeng.xu@kaust.edu.sa }
}

\date[June 18]{OneMath World School \\
Brin Mathematical Research Center\\
University of Maryland\\
June 18, 2026}

\begin{document}

\begin{frame}
	\titlepage
	{
		\scriptsize
%		\vspace{-0.5cm}
%		\begin{center}
%			\large     Joint work with Prof. Jinchao Xu, Dr. Jongho Park
%		\end{center}
%\vspace{-0.5cm}
%			\centerline{Acknowledgement: }
%	\centerline{ KAUST Baseline Research Fund}
		%\begin{figure}
		%\includegraphics[width=0.08\textwidth, height=0.1\textheight]{figures/NSF_4-Color_bitmap_Logo.png}
		%\end{figure}
		
		%\centerline{{\color{blue}NSF} DMS-2111387, KAUST Baseline Fund}
%		\centerline{ NSF DMS-2111387, KAUST Baseline Research Fund}
	}
\end{frame}

\begin{frame}{Outline}
	\tableofcontents
\end{frame}

\section{Introduction}

\begin{frame}{Physics-Informed Neural Networks (PINNs)}
	\begin{itemize}
		\item Approximate PDE solutions with a neural network $u_\theta$
		\item Enforce the governing equation at \textbf{collocation points} via automatic differentiation
		\item Impose \textbf{homogeneous Dirichlet} boundary conditions exactly using a hard constraint (distance function)
		\item Train by minimizing the PDE residual; evaluate against a known exact solution
	\end{itemize}
\end{frame}

\section{One-Dimensional Elliptic Problem}

\begin{frame}{1D Problem Setting (\texttt{pinn.py})}
	\textbf{Domain:} $x \in (0,1)$

	\vspace{0.4em}
	\textbf{PDE:}
	\[
		-u''(x) = \sin(\omega x), \qquad \omega = 10,
	\]
	equivalently $u''(x) + b(x) = 0$ with source $b(x) = \sin(\omega x)$.

	\vspace{0.4em}
	\textbf{Boundary conditions:} homogeneous Dirichlet
	\[
		u(0) = u(1) = 0.
	\]

	\vspace{0.4em}
	\textbf{Exact solution:}
	\[
		u_{\mathrm{exact}}(x)
		= \frac{\sin(\omega x)}{100} - \frac{x\sin(\omega)}{100}.
	\]
\end{frame}

\begin{frame}{1D PINN Formulation (\texttt{pinn.py})}
	\textbf{Hard boundary constraint:}
	\[
		u_\theta(x) = x(1-x)\,\widehat{u}_\theta(x),
	\]
	so $u_\theta(0)=u_\theta(1)=0$ automatically.

	\vspace{0.4em}
	\textbf{Network:} $\widehat{u}_\theta \colon \mathbb{R}\to\mathbb{R}$ is a 2-layer MLP,
	\[
		1 \xrightarrow{\;\tanh\;} 100 \xrightarrow{\;\mathrm{linear}\;} 1,
	\]
	trained in \texttt{float64}.

	\vspace{0.4em}
	\textbf{Loss:} mean squared PDE residual at $500$ interior collocation points,
	\[
		\mathcal{L}(\theta)
		= \frac{1}{N}\sum_{i=1}^{N}
		\bigl(u_\theta''(x_i) + \sin(\omega x_i)\bigr)^2.
	\]

	\vspace{0.4em}
	\textbf{Optimization:} Adam ($6000$ epochs, $\mathrm{lr}=10^{-3}$), then L-BFGS ($100$ steps).
	\textbf{Remark:} In lower dimension, you can use a uniform grid as the collocation points, while in higher dimension we have to use (quasi) Monte-Carlo points. 
\end{frame}

\section{Two-Dimensional Elliptic Problem}

\begin{frame}{2D Problem Setting (\texttt{pinn2d.py})}
	\textbf{Domain:} $(x,y) \in (0,1)^2$

	\vspace{0.4em}
	\textbf{PDE:}
	\[
		-\Delta u(x,y) = f(x,y), \qquad \omega = 10,
	\]
	with Laplacian $\Delta u = u_{xx} + u_{yy}$.

	\vspace{0.4em}
	\textbf{Boundary conditions:} homogeneous Dirichlet on $\partial(0,1)^2$,
	\[
		u = 0 \quad \text{on } x=0,1 \text{ or } y=0,1.
	\]

	\vspace{0.4em}
	\textbf{Manufactured solution:}
	\[
		u_{\mathrm{exact}}(x,y)
		= \frac{\sin(\omega x)\sin(\omega y)}{100},
	\]
	\[
		f(x,y)
		= \frac{2\omega^2\sin(\omega x)\sin(\omega y)}{100}.
	\]
\end{frame}

\begin{frame}{2D PINN Formulation (\texttt{pinn2d.py})}
	\textbf{Hard boundary constraint:}
	\[
		u_\theta(x,y) = x(1-x)\,y(1-y)\,\widehat{u}_\theta(x,y).
	\]

	\vspace{0.4em}
	\textbf{Network:} $\widehat{u}_\theta \colon \mathbb{R}^2\to\mathbb{R}$ is a 2-layer MLP,
	\[
		2 \xrightarrow{\;\tanh\;} 100 \xrightarrow{\;\mathrm{linear}\;} 1.
	\]

	\vspace{0.4em}
	\textbf{Loss:} mean squared residual on an interior tensor-product grid
	($50\times 50$ collocation points, excluding the boundary),
	\[
		\mathcal{L}(\theta)
		= \frac{1}{N}\sum_{i=1}^{N}
		\bigl(\Delta u_\theta(x_i,y_i) + f(x_i,y_i)\bigr)^2.
	\]

	\vspace{0.4em}
	\textbf{Optimization \& evaluation:} same Adam + L-BFGS schedule as the 1D case;
	relative $L^2$ error measured on a finer $100\times 100$ interior test grid.
\end{frame}

\begin{frame}{Remarks}
To understand the performance of the PINNs, it is important to understand how the following factors affect the convergence:  
\begin{itemize}
	\item Change the network architecture, e.g., use different activation functions, different number of layers, different number of neurons, etc.
	\item Change the target function to higher frequency to explore frequency principle. 
	\item Change the number of collocation points to explore the convergence rate.
\end{itemize}
\vfill 
\begin{itemize}
	\scriptsize
	\item Ma, Y., Xu, Z. Q. J., \& Zhang, J. (2021). Frequency principle in deep learning beyond gradient-descent-based training. arXiv preprint arXiv:2101.00747.
	\item Hong, Q., Siegel, J. W., Tan, Q., \& Xu, J. (2022). On the activation function dependence of the spectral bias of neural networks. arXiv preprint arXiv:2208.04924.
\end{itemize}
\end{frame}

\section{Summary}
\begin{frame}{Summary}
	\begin{itemize}
		\item Both codes solve \textbf{stationary elliptic} problems with manufactured solutions
		\item Boundary conditions are enforced through the hard constraint. 
		\item PDE residuals are computed with \textbf{automatic differentiation} (second derivatives in 1D; Laplacian in 2D)
	\end{itemize}
	\vspace{0.4cm}
	\textbf{Question:} Do we have other training methods to solve PDEs using neural networks that can give us better convergence?
	\vfill 
	You are encouraged to explore the following papers to gain more insights into training neural networks for solving PDEs: 
	\begin{itemize}
		\scriptsize 
		\item Mao, T., Xu, J., \& Xu, X. (2026). Solving High-Dimensional PDEs Using Linearized Neural Networks. arXiv preprint arXiv:2601.11771.
		\item Xu, J., \& Xu, X. (2025). Randomized greedy algorithms for neural network optimization in solving partial differential equations. Journal of Scientific Computing, 105(1), 26.
		\item Siegel, J. W., Hong, Q., Jin, X., Hao, W., \& Xu, J. (2023). Greedy training algorithms for neural networks and applications to PDEs. Journal of Computational Physics, 484, 112084.
		\item Park, J., Xu, J., \& Xu, X. (2022). A neuron-wise subspace correction method for the finite neuron method. arXiv preprint arXiv:2211.12031.
	\end{itemize} 
\end{frame}

\begin{frame}
	\centering
	\Large Thank you!
\end{frame}

\end{document}
